% Generated draft: requires mathematical and source review.
\documentclass[11pt]{article}
\usepackage[margin=25mm]{geometry}
\usepackage{fontspec}
\setmainfont{Latin Modern Roman}
\usepackage{amsmath,amssymb,mathtools,unicode-math}
\setmathfont{Latin Modern Math}
\usepackage{enumitem,needspace,xurl,hyperref}
\setlength{\parindent}{0pt}
\setlength{\parskip}{5pt}
\setlength{\emergencystretch}{3em}
\newcommand{\N}{\mathbb{N}}
\newcommand{\Z}{\mathbb{Z}}
\newcommand{\Q}{\mathbb{Q}}
\newcommand{\R}{\mathbb{R}}
\newcommand{\C}{\mathbb{C}}
\newcommand{\F}{\mathbb{F}}
\newcommand{\A}{\mathbb{A}}
\newcommand{\PP}{\mathbb{P}}
\newcommand{\E}{\mathbb{E}}
\newcommand{\eps}{\varepsilon}
\newcommand{\dd}{\,\mathrm{d}}
\begin{document}
\section*{UnsolvedMath ID: 11000231}
\textbf{Source problem code:} catalogue entry
\subsection{Definitions and mathematical statement}
\emph{Problem 3.3.} Given \(\varphi_1,\ldots,\varphi_n\), does there exist an algorithm producing the integer \(M\) asserted in Theorem 3.1? The theorem, the nature of the \(\varphi_i\), their admissibility hypotheses, and their effective input encoding are absent from the supplied record.

Writing \(\mathcal D\) for the intended class of admissible finite tuples and \(R(\varphi_1,\ldots,\varphi_n;M)\) for the unprovided theorem's conclusion, the algorithmic target is a procedure that halts on every input in \(\mathcal D\) and returns an integer \(M\in\Z\) for which that conclusion holds. These symbols expose missing data; they do not define a substitute theorem. If Theorem 3.1 restricts \(M\) to positive integers or requires a threshold valid for all larger exponents, those conditions must be recovered rather than guessed.

The source asks for an effective witness from the tuple, not a proof of bare existence, the least possible integer, or a prescribed complexity bound. It does not state whether \(n\) and the ambient group or surface are fixed or part of the input. No mapping-class, pseudo-Anosov, or independence assumption is silently supplied.
\subsection{Short English statement}
Make the integer in Theorem 3.1 effective from the finite input tuple. The theorem and its hypotheses are missing, so neither the output property nor the admissible inputs can be specified further.
\subsection{Sources}
\begin{itemize}
\item Local UnsolvedMath registry, permanent ID 11000231, source code catalogue entry, Problem 3.3, PDF page 270. No individual proposer is supplied.
\item Farb, \emph{Problems on Mapping Class Groups and Related Topics}. Supplied URL: \url{https://www.math.uchicago.edu/~farb/papers/mcgbook.pdf}; recorded access 2026-07-29; extraction pdf-font-structure. The volume contains 329 explicitly typeset headings across chapters, with repeated labels disambiguated by page. The essential cross-reference Theorem 3.1 is not supplied.
\item Supplied status open and OPEN-TRIAGE. The research summary says only ``Cannot verify.'' Current primary-literature status and redistribution terms for the publicly accessible author/publisher copy remain NEEDS\_REVIEW.
\end{itemize}
\end{document}
